\documentclass[tikz,border=10pt]{standalone}
\usepackage{fontspec}
\usepackage[mathrm=sym]{unicode-math}
\setmainfont{texgyrepagella-regular.otf}[
    BoldFont = texgyrepagella-bold.otf,
    ItalicFont = texgyrepagella-italic.otf,
    BoldItalicFont = texgyrepagella-bolditalic.otf
]
\setmathfont{texgyrepagella-math.otf}[AutoFakeBold = 2.5]
\usepackage{tikz}
\usetikzlibrary{arrows.meta, decorations.markings}

\begin{document}
\begin{tikzpicture}[>=Stealth, line cap=butt, line join=round,
  arrowed/.style={
    decoration={
      markings,
      mark=at position 0.55 with {\arrow{Stealth[length=5pt, width=3.5pt]}}
    },
    postaction={decorate}
  },
  arrowed_double/.style={
    decoration={
      markings,
      mark=at position 0.55 with {\arrow{Stealth[length=5.5pt, width=4.8pt]}}
    },
    postaction={decorate}
  }
]

  % ========================================================
  % ROW 1: (a) Dyson equation with pairing mode
  % ========================================================
  \begin{scope}[shift={(0, 1.4)}]

    % Term 1: Full Green function G (solid line)
    \draw[thick, line width=1.1pt, arrowed] (-4.6, 0) -- (-3.2, 0);

    % Equals sign
    \node[font=\large] at (-2.4, 0) {$=$};

    % Term 2: Free Green function G0 (dashed line)
    \draw[thick, dashed, line width=1pt, arrowed] (-1.6, 0) -- (-0.2, 0);

    % Plus sign
    \node[font=\large] at (0.6, 0) {$+$};

    % Term 3: G0 Sigma G
    \begin{scope}[shift={(1.4, 0)}]
      % Incoming G0
      \draw[thick, dashed, line width=1pt, arrowed] (0, 0) -- (1.4, 0);

      % Bubble
      \coordinate (BL) at (1.4, 0);
      \coordinate (BR) at (3.0, 0);

      % Top arc: double line pairing mode
      \draw[thick, line width=0.9pt, double, double distance=1.6pt, arrowed_double] (BL) to[out=60, in=120] (BR);

      % Bottom arc: dashed line G0
      \draw[thick, dashed, line width=1pt, arrowed] (BL) to[out=-60, in=-120] (BR);

      % Outgoing G
      \draw[thick, line width=1.1pt, arrowed] (BR) -- (4.4, 0);

      % Vertices
      \filldraw[black] (BL) circle (2pt);
      \filldraw[black] (BR) circle (2pt);
    \end{scope}

  \end{scope}

  % ========================================================
  % ROW 2: (b) Gap equation / Pairing mode
  % ========================================================
  \begin{scope}[shift={(0, -1.2)}]

    % LHS: Vertex Delta with pairing line
    \coordinate (D1) at (-4.6, 0);
    \draw[thick, line width=0.9pt, double, double distance=1.6pt, arrowed_double] (D1) -- (-3.0, 0);
    \filldraw[black] (D1) circle (2pt);

    % Equals sign
    \node[font=\large] at (-2.2, 0) {$=$};

    % Middle: Bubble with outgoing pairing mode
    \begin{scope}[shift={(-1.4, 0)}]
      \coordinate (MBL) at (0, 0);
      \coordinate (MBR) at (1.6, 0);

      % Top arc: particle line (solid)
      \draw[thick, line width=1.1pt, arrowed] (MBL) to[out=60, in=120] (MBR);

      % Bottom arc: hole line (dashed, arrow pointing left)
      \draw[thick, dashed, line width=1pt, arrowed] (MBR) to[out=-120, in=-60] (MBL);

      % Outgoing pairing mode
      \draw[thick, line width=0.9pt, double, double distance=1.6pt, arrowed_double] (MBR) -- (2.9, 0);

      % Vertices
      \filldraw[black] (MBL) circle (2pt);
      \filldraw[black] (MBR) circle (2pt);
    \end{scope}

    % "or" text
    \node[font=\normalsize] at (2.6, 0) {$\text{or}$};

    % RHS: Propagator equals amputated vertex
    \begin{scope}[shift={(3.6, 0)}]
      % Pairing mode propagator
      \draw[thick, line width=0.9pt, double, double distance=1.6pt, arrowed_double] (0, 0) -- (1.4, 0);

      % Equals sign
      \node[font=\large] at (2.0, 0) {$=$};

      % Amputated vertex
      \coordinate (AV) at (3.6, 0);
      % Incoming particle (solid)
      \draw[thick, line width=1.1pt, arrowed] (2.7, 0.45) to[out=0, in=135] (AV);
      % Outgoing hole (dashed, pointing away from AV)
      \draw[thick, dashed, line width=1pt, arrowed] (AV) to[out=-135, in=0] (2.7, -0.45);

      % Outgoing pairing mode
      \draw[thick, line width=0.9pt, double, double distance=1.6pt, arrowed_double] (AV) -- (4.8, 0);

      \filldraw[black] (AV) circle (2pt);
    \end{scope}

  \end{scope}

\end{tikzpicture}
\end{document}
